\chapter{Introduction}
\label{ch:introduction}

\section{Motivation}
\label{sec:motivation}

Modern software systems increasingly rely on microservice architectures, in
which business capabilities are decomposed into independently deployable
services connected through network APIs \cite{newman2015building,
dragoni2017microservices}. This architecture improves scalability and
development velocity, but it also makes failures harder to observe and diagnose.
A latency spike in one service may propagate through several downstream calls,
while the operational evidence is split across metrics, logs, traces, and
deployment metadata \cite{sigelman2010dapper, lin2018aiops}.

Traditional centralized monitoring is poorly suited to this environment. Raw
telemetry streams are large, geographically dispersed, and often sensitive. A
monitoring system that centralizes all observations may introduce bandwidth
bottlenecks, latency, and privacy exposure. At the same time, simple static
thresholds are brittle under bursty traffic and non-stationary workloads. These
conditions motivate a monitoring framework that performs inference close to the
telemetry source, learns collaboratively across nodes, and explains anomalies in
terms of service dependencies.

\section{Problem Statement}
\label{sec:problem_statement}

This thesis addresses the following problem: how can a distributed microservice
platform detect anomalies accurately, preserve the privacy of local telemetry,
identify likely root causes, and remain operational under partial failures?
More precisely, the system receives a continuous stream of service-level and
host-level telemetry from multiple nodes, produces an anomaly decision for each
local time window, and attaches operational context to the alert when several
services degrade together. The desired output is not only a binary anomaly
label, but also a reproducible anomaly score, a threshold value, the model
version that produced the decision, and a ranked list of likely root-cause
services when the causal graph contains enough evidence.

The problem is challenging because the monitoring pipeline must balance four
competing requirements:

\begin{itemize}
    \item \textbf{Detection quality}: the model must capture temporal and
    multivariate patterns rather than isolated threshold violations.
    \item \textbf{Privacy}: local telemetry should not be centralized, and
    model updates should reveal as little as possible about individual samples.
    \item \textbf{Operational usefulness}: anomaly alerts should be linked to
    probable root causes and should avoid unnecessary alert fatigue.
    \item \textbf{Deployability}: the system must tolerate node dropout,
    bounded bandwidth, and coordinator restarts.
\end{itemize}

These requirements create trade-offs: centralizing telemetry improves training
visibility but exposes sensitive data; stronger DP protects privacy but can
corrupt small gradients after convergence; lower thresholds increase recall but
cause alert fatigue. The thesis treats anomaly detection as a systems problem.

The scope is bounded: the framework assumes numerical telemetry sampled at
regular intervals and service dependencies recoverable from traces or
deployment metadata. It does not attempt Byzantine-robust aggregation or
complete privacy guarantees, but evaluates whether a practical
edge/federated architecture can improve detection, reduce data movement, and
provide useful diagnostics.

\section{Research Objectives}
\label{sec:research_objectives}

The objective of this work is to design and evaluate a privacy-preserving,
federated anomaly detection framework for distributed microservice systems. The
research is organized around five concrete objectives:

\begin{enumerate}
    \item design an edge-deployable LSTM-Autoencoder (LSTM-AE) for local
    multivariate anomaly detection \cite{malhotra2015lstm};
    \item train the model collaboratively through Federated Averaging (FedAvg)
    while applying client-side differential privacy \cite{mcmahan2017communication,
    dwork2006differential, abadi2016deep};
    \item reduce federated communication overhead through binary serialization,
    Zstandard, LZ4, and adaptive compression;
    \item enrich anomaly alerts with causal-graph-based root cause analysis;
    \item evaluate adaptive Q-learning thresholding and identify the conditions
    under which it reduces or fails to reduce alert fatigue.
\end{enumerate}

These objectives are mapped to the thesis evaluation as follows. The first
objective is tested through the SMD anomaly-detection benchmark and the
LSTM-AE hyperparameter sweep. The second objective is tested through the
FedAvg baseline and the differential privacy experiments, where the updated
results distinguish between the default clipping configuration
($C=1.0$) and the tighter clipping regime ($C=0.01$). The third objective is
tested through payload serialization, Zstandard, LZ4, and adaptive compression
measurements. The fourth objective is tested through causal-graph RCA on
Train-Ticket and Alibaba-scale topologies. The fifth objective is tested
    through the Q-learning thresholding benchmark (10\,020-minute SMD stream,
    seed~42), which is reported as an exploratory result: the policy eventually
    reduces final alert volume, but its early exploration remains unstable and
    therefore requires reward redesign before production use.

The central research questions are therefore:

\begin{enumerate}
    \item Can a federated LSTM-AE match or improve anomaly detection quality
    without centralizing raw telemetry?
    \item How sensitive is the privacy--utility trade-off to the DP noise
    multiplier and clipping norm?
    \item Does communication compression materially reduce federated update
    overhead without changing model semantics?
    \item Can trace-derived causal graphs provide useful root-cause rankings
    for multi-service incidents?
    \item Under what conditions does adaptive thresholding reduce alert fatigue,
    and when does it fail?
\end{enumerate}

\section{Research Contributions}
\label{sec:contributions}

The thesis contributes a complete framework and an empirical evaluation with the
following outcomes:

\begin{itemize}
    \item a federated LSTM-AE that achieves competitive detection performance
    on the Server Machine Dataset (SMD), outperforming the centralized LSTM-AE
    baseline (full results in~\S\ref{sec:eval_detection});
    \item a differential privacy benchmark showing a sharp utility transition:
    mild noise ($\sigma \leq 0.0005$, $\varepsilon \geq 10$) preserves
    near-baseline performance, while stronger noise requires careful clipping
    (\S\ref{subsec:eval_dp_impact});
    \item a clipping sensitivity analysis showing that tight clipping recovers
    detection utility under moderate noise, while also exposing the practical
    privacy tension created by very small effective noise
    (\S\ref{subsec:eval_clipping});
    \item a causal graph RCA engine evaluated on Train-Ticket and
    Alibaba-derived topologies (\S\ref{sec:eval_rca});
    \item a Q-learning thresholding study that reduces final alert volume while
    motivating warm starts and reward shaping due to unstable early exploration
    (\S\ref{sec:eval_threshold});
    \item a communication and scalability evaluation showing that binary
    payloads are substantially smaller than text payloads and that FedAvg
    aggregation scales approximately linearly to 1\,000 nodes
    (\S\ref{sec:eval_fl_comm}, \S\ref{sec:eval_scalability}).
\end{itemize}

The novelty of this thesis is therefore system-level rather than component-only.
Prior work commonly studies deep anomaly detection, federated optimization,
differential privacy, RCA, or adaptive thresholds as separate mechanisms. This
work evaluates how these mechanisms behave when they are connected in one
edge-deployable monitoring architecture.

This integration is important because the interfaces change the result:
threshold calibration changes the usefulness of reconstruction errors, DP
clipping changes the utility of federated updates, compression changes
deployment feasibility, and RCA changes whether an alert is actionable. The
strongest empirical result is the federated LSTM-AE detection performance,
while the main cautionary results are the DP clipping tension and the
Q-learning cold-start failure. Reporting both positive and negative outcomes is
part of the contribution, because it makes the framework reproducible and gives
clear directions for future improvement.
